\documentclass{ximera}
\title{Systems of Linear Equations in Three Variables}

\outcome{Solve systems of three linear equations in three variables by elimination.}
\outcome{Classify systems as consistent, inconsistent, or dependent.}
\outcome{Set up and solve three-variable systems from word problems.}

\begin{document}
\begin{abstract}
\end{abstract}
\maketitle

\textbf{Learning Objectives:}
\begin{itemize}
  \item Solve systems of three linear equations in three variables by elimination.
  \item Classify systems as consistent, inconsistent, or dependent.
  \item Set up and solve three-variable systems from word problems.
\end{itemize}
\medskip

\section*{Quick Check}

\begin{problem}
Which of the following systems has the solution $(x,y,z) = (1,2,-1)$?
\begin{hint}
Substitute the values into each equation in each system and check which one is satisfied by all three equations simultaneously.
\end{hint}
\begin{multipleChoice}
  \choice{$\begin{cases} x+y+z=2 \\ 2x-y+3z=-1 \\ x-2y+z=-2 \end{cases}$}
  \choice[correct]{$\begin{cases} x+y+z=2 \\ 2x-y+3z=-3 \\ x-2y+z=-4 \end{cases}$}
  \choice{$\begin{cases} x+y+z=1 \\ 2x-y+3z=-3 \\ x-2y+z=-2 \end{cases}$}
  \choice{$\begin{cases} x+y+z=2 \\ 2x-y+3z=-3 \\ x-2y+z=0 \end{cases}$}
\end{multipleChoice}
\end{problem}

\begin{problem}
Which of the following systems does \textbf{not} have a solution?
\begin{hint}
An inconsistent system produces a contradiction (e.g., $0 = 5$) when you attempt to eliminate variables.
\end{hint}
\begin{multipleChoice}
  \choice{$\begin{cases} x+y+z=6 \\ 2x+2y+2z=12 \\ x-y+z=0 \end{cases}$}
  \choice{$\begin{cases} x+y+z=6 \\ x-y+z=0 \\ 2x+2z=8 \end{cases}$}
  \choice{$\begin{cases} x+y+z=6 \\ 2x+2y+2z=10 \\ x-y+z=0 \end{cases}$}
  \choice[correct]{$\begin{cases} x+y+z=6 \\ x-y+z=0 \\ y=1 \end{cases}$}
\end{multipleChoice}
\end{problem}

\begin{problem}
A nut mix is made from almonds ($a$ \$/lb), cashews ($c$ \$/lb), and peanuts ($p$ \$/lb). Order 1: 3 lb almonds, 2 lb cashews, 5 lb peanuts for \$24. Order 2: 1 lb almonds, 4 lb cashews, 2 lb peanuts for \$16. Order 3: 2 lb almonds, 1 lb cashew, 3 lb peanuts for \$13. Which system correctly models this situation?
\begin{hint}
Each order gives one equation: (price per lb) $\times$ (quantity) = total cost.
\end{hint}
\begin{multipleChoice}
  \choice[correct]{$\begin{cases} 3a+2c+5p=24 \\ a+4c+2p=16 \\ 2a+c+3p=13 \end{cases}$}
  \choice{$\begin{cases} 3a+2c+5p=16 \\ a+4c+2p=24 \\ 2a+c+3p=13 \end{cases}$}
  \choice{$\begin{cases} a+c+p=24 \\ a+c+p=16 \\ a+c+p=13 \end{cases}$}
  \choice{$\begin{cases} 3a+2c+5p=13 \\ a+4c+2p=24 \\ 2a+c+3p=16 \end{cases}$}
\end{multipleChoice}
\end{problem}

\begin{problem}
An animal shelter houses dogs ($d$), cats ($c$), and rabbits ($r$). There are 30 animals total. Dogs eat 3 lb/day, cats eat 1 lb/day, and rabbits eat 0.5 lb/day, for a total of 90 lb/day. The number of cats is twice the number of dogs. Which system represents this situation?
\begin{hint}
Write one equation for the total number of animals, one for the total food consumed, and one for the relationship between cats and dogs.
\end{hint}
\begin{multipleChoice}
  \choice[correct]{$\begin{cases} d+c+r=30 \\ 3d+c+0.5r=90 \\ c=2d \end{cases}$}
  \choice{$\begin{cases} d+c+r=90 \\ 3d+c+0.5r=30 \\ c=2d \end{cases}$}
  \choice{$\begin{cases} d+c+r=30 \\ 3d+c+0.5r=90 \\ c=d/2 \end{cases}$}
  \choice{$\begin{cases} d+c+r=30 \\ 3d+2c+0.5r=90 \\ c=2d \end{cases}$}
\end{multipleChoice}
\end{problem}

\begin{problem}
A nutritionist designs a meal using foods A, B, and C. The total number of servings is 10. Food A contributes 2 g protein/serving, B contributes 4 g, and C contributes 3 g, for a total of 30 g. Each serving of A, B, C costs \$1, \$2, \$3 respectively, for a total cost of \$20. Which system correctly models this situation?
\begin{hint}
Write one equation for servings, one for protein, and one for cost.
\end{hint}
\begin{multipleChoice}
  \choice[correct]{$\begin{cases} a+b+c=10 \\ 2a+4b+3c=30 \\ a+2b+3c=20 \end{cases}$}
  \choice{$\begin{cases} 2a+4b+3c=10 \\ a+b+c=30 \\ a+2b+3c=20 \end{cases}$}
  \choice{$\begin{cases} a+b+c=30 \\ 2a+4b+3c=20 \\ a+2b+3c=10 \end{cases}$}
  \choice{$\begin{cases} a+b+c=10 \\ 2a+4b+3c=20 \\ a+2b+3c=30 \end{cases}$}
\end{multipleChoice}
\end{problem}


\bigskip
\section*{Extended Practice}

\begin{problem}
Solve the system of equations. If the system does not have a unique solution, state how many solutions it has.
\begin{enumerate}
  \item[(a)] $\ds\begin{cases} 5x = 2y-3z-3\\ 4y = -1-10x-5z\\ 2z = 5x - 6y \end{cases}$
  \medskip
  \item[(b)] $\ds\begin{cases} 2x + 5z = 2\\ 3y - 7z = 9\\ -5x+9y = 22 \end{cases}$
  \medskip
  \item[(c)] $\ds\begin{cases} -x+4y+2z = 4\\ x - 3y - z = -2\\ -x+y-z = -2 \end{cases}$
\end{enumerate}

\begin{solution}
\begin{enumerate}
  \item[(a)] Put in standard form:
  $$\begin{cases} 5x-2y+3z=-3 & (E1)\\ 10x+4y+5z=-1 & (E2)\\ -5x+6y+2z=0 & (E3) \end{cases}$$
  Eliminate $x$: Add $(E1)+(E3)$: $4y+5z=-3$ $(E4)$. Add $(E2)+2(E3)$: $16y+9z=-1$ $(E5)$.\\
  Eliminate $y$ from $(E4),(E5)$: Add $-4(E4)+(E5)$: $-11z=11 \implies z=-1$.\\
  From $(E4)$: $4y-5=-3 \implies y=\dfrac{1}{2}$.\\
  From $(E3)$: $-5x+3-2=0 \implies x=\dfrac{1}{5}$.\\
  Solution: $\left(\dfrac{1}{5},\dfrac{1}{2},-1\right)$.

  \item[(b)] Note each equation is missing one variable. Eliminate $y$: Add $-3(E2)+(E3)$: $-5x+21z=-5$ $(E4)$.\\
  Now solve $(E1),(E4)$: eliminate $x$ by adding $5(E1)+2(E4)$: $10x+25z=10$ and $-10x+42z=-10$ gives $67z=0 \implies z=0$.\\
  From $(E1)$: $2x=2 \implies x=1$. From $(E2)$: $3y=9 \implies y=3$.\\
  Solution: $(1,3,0)$.

  \item[(c)] Eliminate $x$: Add $(E1)+(E2)$: $y+z=2$ $(E4)$. Add $(E2)+(E3)$: $-2y-2z=-4$ $(E5)$.\\
  Note $(E5)=-2(E4)$. Adding $2(E4)+(E5)$: $0=0$.\\
  \textbf{Dependent system} — infinitely many solutions. Let $z=t$; then $y=2-t$ and from $(E3)$: $x=4-2t$.\\
  General solution: $(4-2t,\;2-t,\;t)$ for any real $t$.
\end{enumerate}
\end{solution}
\end{problem}


\begin{problem}
Set up and solve a system of linear equations for each scenario. Clearly define your variables.
\begin{enumerate}
  \item[(a)] Ben inherited \$120,000 and invested it in three accounts: a money market account (2.2\% simple interest), a stock (6\% return), and a mutual fund ($-2\%$ return). He invested \$10,000 more in the stock than in the mutual fund, and his net gain for the year was \$2,820. Find the amount in each account.
  \medskip
  \item[(b)] The largest angle of a triangle is $18^\circ$ more than the sum of the other two angles. The smallest angle is half the middle angle. Find all three angle measures.
  \medskip
  \item[(c)] A theater sells Section A tickets for \$50, Section B for \$30, and Section C for \$20. For one play, 4000 tickets were sold for total revenue of \$120,000. Section B sold 1000 more tickets than the other two sections combined. How many tickets were sold in each section?
\end{enumerate}

\begin{solution}
\begin{enumerate}
  \item[(a)] Let $m, s, f$ be amounts in money market, stock, and mutual fund respectively.
  $$\begin{cases} m+s+f=120{,}000 \\ 0.022m+0.06s-0.02f=2{,}820 \\ s-f=10{,}000 \end{cases}$$
  Multiply $(E2)$ by 1000, eliminate $m$ by adding $-22(E1)+(E2)$: $38s-42f=180{,}000$ $(E4)$.\\
  Eliminate $s$ from $(E3),(E4)$: add $-38(E3)+(E4)$: $-4f=-200{,}000 \implies f=50{,}000$.\\
  From $(E3)$: $s=60{,}000$. From $(E1)$: $m=10{,}000$.\\
  \textbf{Ben invested \$10,000 in the money market, \$60,000 in stock, and \$50,000 in the mutual fund.}

  \item[(b)] Let $x, y, z$ be the largest, middle, and smallest angles respectively.
  $$\begin{cases} x+y+z=180 \\ x=y+z+18 \\ z=\frac{1}{2}y \end{cases}$$
  Standard form: $\begin{cases} x+y+z=180 \\ x-y-z=18 \\ y-2z=0 \end{cases}$. Add $(E1)+(-E2)$: $2y+2z=162$ $(E4)$.\\
  Add $(E3)+(E4)$: $3y=162 \implies y=54$. Then $z=27$, $x=99$.\\
  \textbf{The angles are $99^\circ$, $54^\circ$, and $27^\circ$.}

  \item[(c)] Let $a, b, c$ be tickets sold in Sections A, B, C respectively.
  $$\begin{cases} a+b+c=4{,}000 \\ 50a+30b+20c=120{,}000 \\ b=a+c+1{,}000 \end{cases}$$
  Standard form $(E3)$: $-a+b-c=1{,}000$. Add $(E1)+(E3)$: $2b=5{,}000 \implies b=2{,}500$.\\
  Eliminate $a$ from $(E1),(E2)$: add $-50(E1)+(E2)$: $-20b-30c=-80{,}000$.\\
  With $b=2{,}500$: $-30c=-30{,}000 \implies c=1{,}000$. Then $a=500$.\\
  \textbf{500 Section A tickets, 2,500 Section B tickets, 1,000 Section C tickets.}
\end{enumerate}
\end{solution}
\end{problem}

\end{document}
